\chapter{Introdução}

Projetar e fiscalizar obras na engenharia civil brasileira ainda envolve uma quantidade enorme de retrabalho. Parte desse desperdício vem da forma como os projetos são representados: o desenho 2D, padrão histórico do setor, esconde colisões e omite informações que só aparecem no canteiro. A migração para modelos tridimensionais e, mais especificamente, para a metodologia \textit{Building Information Modeling} ou BIM vem como resposta a esse problema \cite{eastman2011}.

O BIM organiza o planejamento e a execução da obra em torno de um modelo 3D detalhado. A metodologia divide-se em dimensões que vão do 3D ao 7D: 3D para o modelo geométrico em si; 4D para a programação temporal; 5D para custos; 6D para sustentabilidade; e 7D para o gerenciamento da operação pós-construção \cite{eastman2011}. BIM não é um \textit{software}, e sim um conceito que se apoia em sistemas informatizados como o CAD. Quem implementa esse conceito são ferramentas como o Revit, da Autodesk, e o ArchiCAD, da Graphisoft.

Modelar em BIM permite visualizar interferências antes da obra começar, o que reduz o retrabalho em campo. As aplicações típicas vão da análise de projetos ao orçamento, passando pelo planejamento, pela coordenação modular e pela gestão do ciclo de vida \cite{eastman2011}.

O BIM, contudo, não resolve sozinho a verificação dos projetos contra as normas regulatórias. As normas regulatórias (ou normas técnicas) são documentos que estabelecem requisitos, critérios e procedimentos obrigatórios para assegurar a segurança, a qualidade e o desempenho das edificações; no Brasil, são editadas pela Associação Brasileira de Normas Técnicas (ABNT) na forma das NBRs, como a NBR~9050, voltada à acessibilidade. Cada jurisdição mantém o próprio conjunto de regras, em geral redigidas em linguagem natural e de caráter interpretativo, o que faz da verificação automática de conformidade um problema em aberto \cite{dimyadi2013,solihin2015}.

Tornar essas normas verificáveis por máquina exige convertê-las em uma representação computável. Com essa finalidade, \citeonline{hjelseth2011} propuseram a metodologia RASE, que reescreve cada trecho normativo como dado estruturado, decomposto em quatro atributos: \textbf{R}equisito (o que precisa ser atendido), \textbf{A}plicabilidade (onde a regra se aplica), \textbf{S}eleção (opções para cumpri-la) e \textbf{E}xceção (casos em que a regra não vale). Uma vez organizada em RASE, a norma viabiliza ferramentas integradas a \textit{softwares} BIM, como o Revit, capazes de validar o projeto elemento a elemento \cite{santos2021}. O gargalo está antes disso: a conversão da norma para RASE ainda é feita manualmente, regra por regra, e, como um modelo BIM pode reunir centenas ou milhares de elementos sujeitos a normas, o custo desse trabalho cresce rapidamente.

Automatizar essa conversão exige reconhecer padrões em textos redigidos em linguagem natural, tarefa para a qual o Aprendizado de Máquina (AM) se mostra promissor. A motivação é reforçada pelo próprio BIM: quanto maior a sua adoção, maior o volume de dados gerado por cada projeto, o que atrai cientistas de dados, profissionais cuja função é transformar grandes volumes de informação em valor para as organizações \cite{davenport2012}. O AM, subárea da Inteligência Artificial (IA), estuda como construir algoritmos capazes de aprender padrões diretamente dos dados, sem programação explícita das regras \cite{tan2006}; a partir dos padrões aprendidos, esses modelos fazem previsões e respondem a situações novas. Na Engenharia Civil, a IA já automatiza etapas que antes dependiam de inspeção humana, da detecção de patologias em estruturas à previsão de prazos de obra \cite{rane2023,liu2022}.

Três exemplos recentes ilustram esse movimento. \citeonline{liu2022} combinaram BIM, AM e regras de \textit{design} para otimizar o processo de montagem da construção, mostrando ganhos em eficiência e custo. \citeonline{alawadhi2023} treinaram Redes Neurais Artificiais sobre dados BIM para reconhecer objetos de construção em fotos: as imagens sintéticas geradas a partir do modelo produziram segmentação semântica eficaz em fotografias reais, abrindo caminho para o uso de dados sintéticos em problemas arquitetônicos. Em outra frente, \citeonline{saka2024} alimentaram um modelo com 2.280 projetos de demolição para prever a circularidade de resíduos: o sistema estima as quantidades de material reciclável e destina os resíduos para o aterro adequado, apoiando o planejamento.

Quando o objeto de análise é o próprio texto normativo, a automação recai sobre o Processamento de Linguagem Natural (PLN), subárea da IA dedicada ao tratamento computacional da linguagem humana em tarefas como tradução, sumarização e classificação de texto. Dentro do PLN, a IA Generativa (IAG) reúne os modelos capazes de produzir novos conteúdos (texto, imagem ou áudio) a partir do que aprenderam durante o treinamento. Entre suas manifestações mais relevantes estão os Modelos de Linguagem de Grande Porte (\textit{Large Language Models} ou LLMs) \cite{zhao2023}, categoria que inclui o ChatGPT, da OpenAI \cite{openai2024}, e o Llama 3, da Meta \cite{llama3}.

Os LLMs são sistemas de IA voltados a entender e produzir linguagem humana em escala. O treinamento envolve volumes massivos de texto, o que permite ao modelo capturar nuances linguísticas e contextuais. Quase todos os LLMs atuais usam alguma variação de arquitetura \textit{transformer}, escolhida por sua eficiência no processamento de sequências longas \cite{zhao2023}.

A adaptação desses modelos a domínios específicos pode ser feita por diferentes caminhos. O mais leve e acessível deles é a \textit{Engenharia de Prompt} (EP): em vez de retreinar pesos ou montar uma infraestrutura de recuperação, ajusta-se apenas a instrução enviada ao modelo.

No setor AEC, os LLMs podem auxiliar tanto na leitura das normas quanto na sua conversão para formatos processáveis por máquina \cite{fuchs2024}. Só que a literatura ainda tem lacunas: uma delas é justamente a interpretação de leis, códigos e padrões regulatórios, cuja complexidade semântica historicamente desafiou os métodos de IA anteriores \cite{fuchs2024}.

Diante dessa lacuna, este trabalho decompõe a conversão de normas para RASE em três níveis sucessivos de normalização, apoiados unicamente em Engenharia de Prompt: o nível \textbf{N1} segmenta a sentença normativa em regras atômicas; o \textbf{N2} identifica, em cada regra, os operadores RASE; e o \textbf{N3} estrutura o resultado em JSON. Sobre essa decomposição, comparam-se seis LLMs \textit{open-source} (Llama, Dolphin, Gemma, Mistral, Alpaca e Qwen), servidos localmente, em cinco experimentos que avaliam cada nível isoladamente e em \textit{pipeline} encadeado.

A qualidade das saídas é aferida por nove métricas de similaridade textual, distribuídas em quatro famílias: lexicais (FuzzyWuzzy e TF-IDF), semânticas (SBERT, BERTimbau e Multilingual), de distância (\textit{Word Mover's Distance} com FastText e NILC) e voltadas para geração (BERTScore e ROUGE-L). A essas somam-se quatro métricas de classificação (Acurácia, Precisão, Revocação e F1-score).


\section{Objetivos}

\subsection*{Objetivo geral}

Fazer uma análise comparativa do uso de seis Modelos de Linguagem de Grande Porte (LLMs) \textit{open-source} (Llama, Dolphin, Gemma, Mistral, Alpaca e Qwen) na conversão automática de normas técnicas brasileiras da AEC (em especial a NBR 9050) para o formato RASE, utilizando exclusivamente Engenharia de Prompt em três níveis sucessivos de normalização (N1, N2 e N3).

\subsection*{Objetivos específicos}

\begin{itemize}
    \item Definir e implementar os três níveis de normalização (N1, N2 e N3) para a conversão de normas em RASE.
    \item Projetar \textit{prompts} especializados para cada nível.
    \item Comparar os seis LLMs \textit{open-source} nos cinco experimentos propostos.
    \item Avaliar a qualidade das saídas com métricas de similaridade textual e de classificação.
    \item Analisar a propagação de erros ao longo do \textit{pipeline} e o custo computacional por modelo.
\end{itemize}


\section{Metodologia}

Para alcançar os objetivos específicos, a conversão de normas para RASE foi modelada como um \textit{pipeline} de três níveis (N1, N2 e N3), com \textit{prompts} dedicados a cada nível e a cada operador, construídos com as técnicas \textit{Prompt Direct}, \textit{Chain-of-Thought} e \textit{Few-Shot}. Os seis modelos \textit{open-source} (Llama, Dolphin, Gemma, Mistral, Alpaca e Qwen) foram servidos localmente via Ollama, sob protocolo idêntico de \textit{hardware}, \textit{prompts} e \textit{dataset}. O \textit{dataset} reúne 79 normas brasileiras anotadas, com ênfase na NBR~9050 e referência em todos os níveis de normalização.

A avaliação foi organizada em cinco experimentos: \textbf{EN1} (N1 isolado), \textbf{EN2} (N2 com o N1 de referência como entrada), \textbf{EN1N2} (\textit{pipeline} encadeado N1~$\rightarrow$~N2), \textbf{EN3} (N3 com o N2 de referência) e \textbf{EN1N2N3} (\textit{pipeline} completo N1~$\rightarrow$~N2~$\rightarrow$~N3). Em cada experimento, as saídas dos modelos são comparadas, par a par, à referência do \textit{dataset} pelas métricas de similaridade textual e de classificação já apresentadas, complementadas pela medição do tempo de execução por modelo.


\section{Contribuições}

As contribuições deste trabalho situam-se na interseção entre a Arquitetura, Engenharia e Construção (AEC) e a Computação: para a AEC, aplica modelos de PLN à checagem de normas de construção e propõe níveis sucessivos de normalização sobre a metodologia RASE de \citeonline{hjelseth2011}; para a Computação, leva LLMs consolidados na literatura a um domínio inédito e os submete a uma avaliação sistemática de desempenho. O detalhamento das contribuições é retomado na Conclusão.


\section{Organização do Documento}

O restante do documento segue a estrutura abaixo. O \textbf{Capítulo~2: Fundamentação Teórica} cobre os conceitos necessários ao trabalho: BIM, as normas de engenharia (em especial as normas regulatórias brasileiras na AEC) e a metodologia RASE; na sequência, o Aprendizado de Máquina e o PLN; as métricas de similaridade textual adotadas; as redes neurais artificiais (com LSTM, BiLSTM e \textit{Transformers}) e os LLMs; encerrando com os trabalhos relacionados. O \textbf{Capítulo~3: Metodologia} detalha os níveis de normalização N1, N2 e N3, as técnicas de Engenharia de Prompt usadas, os seis LLMs avaliados, o \textit{dataset} de 79 normas, o pipeline de execução e as métricas de validação. Já o \textbf{Capítulo~4: Resultados} traz o desempenho por modelo e por métrica nos cinco experimentos (EN1, EN2, EN1N2, EN3 e EN1N2N3), somado aos tempos de execução e à análise comparativa. Por fim, o \textbf{Capítulo~5: Conclusão} consolida as contribuições, expõe as limitações e aponta direções de trabalho futuro.